\chapter{Security hardening: the current secure-by-construction model}

\noindent\textbf{Learning path: Fast path: mandatory before shipping.} Consolidate the security invariants that must hold across the whole application.\par\medskip

Velran security is not optional middleware: authority is divided across language, compiler, runtime and server. Make security properties compiler invariants where possible; otherwise enforce a secure platform default or fail closed with a stable diagnostic.

\noindent\textbf{First read.} Focus on authority/access control, input/browser boundaries, resource safety, and the production security gate. Use the remaining sections as a hardening reference.\par\medskip

\section{What the secure model does -- and does not -- decide}

Velran can make raw SQL, arbitrary redirect URLs, ambient DB/network authority, unbounded list queries, forged redaction and several other unsafe states unrepresentable. It can prove that a mutation used an authorization proof. It cannot decide whether ``Editor may publish this particular article'' is the correct business rule. That decision remains in the application's permission/object/tenant policy.

A useful review rule is: \emph{mechanism should be compiler/runtime-owned; business intent should remain explicit application code}.

\section{Authority and access control}

Every route has an explicit access policy. There is no implicit public endpoint. Authorization is layered:

\begin{enumerate}
  \item trusted authentication establishes the principal;
  \item route policy determines whether the handler may be entered;
  \item object authorization proves access to a loaded record;
  \item mutation authorization proves that an UPDATE/DELETE key comes from authorized object authority;
  \item named permissions and MFA can be part of a handler contract;
  \item tenant-scoped data requires active-tenant membership proof;
  \item effects/capabilities make DB and outbound-network authority visible rather than ambient.
\end{enumerate}

Validation proof is never authorization proof. A valid \code{ArticleId} still does not prove that the current principal may mutate that article.

\section{Critical operations, idempotency and transaction outcome}

A critical contract combines security obligations rather than scattering them through boilerplate:

\begin{lstlisting}[language=Velran]
critical Payment {
    permission BillingWrite
    mfa
    transaction
    idempotency
    audit PaymentCharged
}
\end{lstlisting}

Idempotency binds the replay key to both request fingerprint and principal scope. Reusing the key with a different payload becomes a conflict instead of a second side effect.

Database commit acknowledgement is also explicit. A captured transaction returns the closed sum type \code{Committed | RolledBack | CommitUnknown}, and \code{match} is exhaustive with no wildcard fallthrough. An idempotent critical transaction must capture and handle the outcome; \code{CommitUnknown} cannot disappear into a generic retryable DB error.

\section{Input, injection and browser boundaries}

The major injection surfaces are structural rather than convention-based:

\begin{itemize}
  \item SQL uses compiler-owned statements and typed binds; string-interpolated raw SQL is rejected;
  \item HTML/template output is structured and sensitivity-aware;
  \item redirects use typed local route calls rather than string URLs;
  \item response headers use typed names and bounded validated values;
  \item filenames, media types and Content-Disposition are typed;
  \item outbound calls use named integrations and relative literal paths, not arbitrary URLs;
  \item generated CSP defaults to least privilege, including \code{connect-src 'none'};
  \item compressed inbound HTTP bodies are rejected until a future bounded decompression surface is explicitly designed.
\end{itemize}

\section{Trust, domain types and disclosure}

Request-boundary values begin untrusted. Domain validation can refine them to validated nominal types, but this does not erase sensitivity or create authorization.

\code{Secret<T>} and \code{Sensitive<T>} metadata follows values through calls. Full database/domain models are not generically serialized; public output uses explicit projections. This prevents a newly added model field from silently appearing in an existing API.

For logs and monitoring, \code{Redacted<T>} cannot be forged with a type annotation. It is produced only by the trusted \code{redact(...)} builtin, which discards the original representation and yields a safe redacted value.

\section{Authentication and session hardening}

The browser session is platform-owned. Login rotates session authority, logout invalidates it, and auth-generation changes revoke older sessions after security-relevant account changes. Cookies use the hardened platform policy rather than application-selected flags.

Online authentication abuse is limited in three dimensions:

\begin{itemize}
  \item source + principal;
  \item principal across sources;
  \item source across principals.
\end{itemize}

MFA/recovery attempts have a stricter stage limit. Limiter-store failure is fail closed, and client-visible credential errors remain deliberately generic to reduce account enumeration.

\section{Credentials, tokens and cryptography}

Purpose-specific types prevent accidental cross-use of passwords, token hashes, session tokens, reset tokens, CSRF tokens and cryptographic keys.

Passwords use platform-owned Argon2id policy. Bearer tokens are persisted as purpose-specific hashes and use expiry-aware verification; reset grants, session revocation and session rotation have atomic lifecycle contracts.

Cryptographic keys have purpose and lifecycle. For user-data encryption the high-level primitive is fixed to AES-256-GCM with a platform-generated nonce and versioned authenticated envelope. The lifecycle is:

\begin{lstlisting}
Active    -> encrypt + decrypt
Retiring  -> decrypt only
Retired   -> neither
\end{lstlisting}

Application code cannot select a weaker cipher or supply its own nonce.

\section{Verified webhooks and safe uploads}

A webhook is not an unauthenticated public POST. The route has separate webhook authority: raw-body signature verification, timestamp/replay-window checking and an atomic replay claim happen before the handler runs.

File upload follows a state machine:

\begin{lstlisting}
multipart bytes
    -> private staged file
    -> byte-authoritative inspection
    -> verified image
    -> handler success
    -> atomic publish
\end{lstlisting}

Client filename and MIME are metadata, not storage or type authority. Public media serving only exposes directories declared by verified image publish routes.

\section{Resource safety is a security invariant}

Velran bounds both individual operations and the whole request:

\begin{itemize}
  \item request/parser hard limits;
  \item route resource profiles under a platform ceiling;
  \item instruction/allocation budgets and hard request deadline;
  \item static row bounds for database list queries plus runtime row caps;
  \item bounded request collections;
  \item upload byte/pixel limits;
  \item bounded outbound request and response bodies;
  \item cumulative external-I/O accounting across request and outbound traffic.
\end{itemize}

General unrestricted collection growth is intentionally not exposed until the same cardinality guarantees can be preserved.

\section{Production configuration is a security contract}

An application can declare strict production requirements:

\begin{lstlisting}[language=Velran]
production {
    https required;
    debug disabled;
    hsts required;
    database tls required;
}
\end{lstlisting}

Startup fails closed when the deployment cannot satisfy the contract. Production mode rejects insecure development cookies, development-mode source reload, missing-Origin exceptions and insecure remote database transport; transactional \code{reload.mode = "rolling"} remains allowed. Reload cannot silently change the deployment security contract; changing it requires restart and operator review.

\section{Supply-chain state is part of the release}

The release records more than crate versions. Three artifacts belong together:

\begin{lstlisting}
Cargo.lock
SUPPLY-CHAIN-CAPABILITIES.txt
VELRAN-SUPPLY-CHAIN.lock
\end{lstlisting}

Direct external dependency edges have explicit security capabilities such as network, filesystem, crypto or parser. Current provenance policy accepts checksummed crates.io packages only. The supply-chain envelope seals both the Cargo lock and capability policy, and release tooling verifies \code{cargo metadata --locked} before sealing or packaging.

\section{Security events, redaction and alerts}

Critical operations already emit typed audit events. The platform also emits structured security events for authentication abuse, MFA failures, access-policy denial, CSRF/origin/CORS rejection, webhook verification failure, idempotency conflict and resource exhaustion.

The event schema does not accept arbitrary request payloads. Principal/source values are redacted in emitted events; bounded in-memory correlation may use raw keys without writing them to logs. Platform-owned burst thresholds and cooldowns turn repeated failures into alerts without allowing alert floods to become a new availability problem.

\section{Production security gate}

Before handing a release to production, verify at least:

\begin{enumerate}
  \item the real toolchain gates pass: format, locked metadata, check, tests and \code{./verify.sh};
  \item production policy is present where required and startup compatibility is green;
  \item HTTPS, HSTS, secure cookie and trusted-proxy topology are correct;
  \item DB runtime credentials are least privilege and remote DB TLS is enforced;
  \item auth/session/rate-limit shared state matches the deployment topology;
  \item tenant, object and mutation authorization negative tests are green;
  \item critical/idempotent operations handle \code{CommitUnknown} explicitly;
  \item webhook verification/replay and upload publish-state tests are green;
  \item outbound integrations resolve only through reviewed named egress policy;
  \item resource/deadline limits are configured and observed;
  \item security-event redaction and burst alerts are verified;
  \item supply-chain and release-manifest state are current and verified.
\end{enumerate}

\section{Review questions for every new feature}

Ask five questions:

\begin{enumerate}
  \item \textbf{What comes from outside?} Which fields, files, identifiers or upstream data are untrusted?
  \item \textbf{Where does it become typed or validated?} Which boundary creates the proof?
  \item \textbf{Who owns authority?} Session, database record, tenant membership, trusted config or named capability?
  \item \textbf{What happens in ambiguous failure?} Is there a fail-closed, rollback, cleanup or explicit unknown state?
  \item \textbf{What bounds the resource?} Bytes, rows, items, time, rate, CPU or cumulative I/O?
\end{enumerate}

If any answer is vague, the feature's security model is unfinished.

\section{Verified companion examples}
Security hardening is easier to review against executable fixtures. The canonical companions here are \path{examples/resource-profile/app.vrn}, \path{examples/rate-limit/app.vrn}, and \path{examples/typed-security-events/main.vrn}.

\section{Practice: map controls to threat paths}
\paragraph{Exercise mode: independent.} Use the independent method defined in the preface.

Choose three concrete threats and trace which independent controls stop or limit each one. Prefer layered answers: typed input plus authorization plus transaction invariants is stronger than a single validation check expected to do everything.

\section{Common mistake: treating hardening as a final deployment phase}
Hardening that depends on application structure must be designed before production. Route policy, capability boundaries, object authorization, bounded inputs, and named egress are architecture, not post-release configuration polish.

\section{Running project checkpoint: Secure Notes}
Review Secure Notes as an attacker would: identifiers, form fields, publication transitions, media references, JSON endpoints, and outbound integration. Document which boundary rejects each invalid path and which event should be observable when the rejection matters operationally.
